% Source printed pp. 93-98.
\subsection{预备: 群胚的纤维积}\label{src:3.4.6-subsection}
预叠的纤维积会显露出许多 $2$-范畴的微妙之处. 理解纤维积对于使用代数叠至关重要, 因而建议读者花时间练习许多例子. 从群胚的纤维积入手很有帮助.
\begin{construction}\label{src:3.4.28}
设 $f:\mathcal C\to\mathcal D$ 与 $g:\mathcal D'\to\mathcal D$ 是群胚的函子. 定义群胚 $\mathcal C\times_{\mathcal D}\mathcal D'$ 为由三元组 $(c,d',\gamma)$ 组成的范畴, 其中 $c\in\mathcal C$, $d'\in\mathcal D'$ 是对象, $\gamma:f(c)\xrightarrow{\sim}g(d')$ 是 $\mathcal D$ 中的同构.
态射 $(c_1,d_1',\gamma_1)\to(c_2,d_2',\gamma_2)$ 是态射 $\chi:c_1\xrightarrow{\sim}c_2$ 与 $\lambda:d_1'\xrightarrow{\sim}d_2'$ 的数据, 使得
\[
\begin{tikzcd}
f(c_1)\arrow[r,"f(\chi)"]\arrow[d,"\gamma_1"']&
f(c_2)\arrow[d,"\gamma_2"]\\
g(d_1')\arrow[r,"g(\lambda)"']&g(d_2')
\end{tikzcd}
\]
交换.
\end{construction}
\begin{exercise}[细节]\label{src:3.4.29}
表述群胚纤维积的泛性质, 并证明 $\mathcal C\times_{\mathcal D}\mathcal D'$ 满足这一性质.
\end{exercise}
下面的习题预示了涉及商叠的重要 Cartesian 图.
\begin{exercise}[重要, 有益的练习]\label{src:3.4.30}
设群 $G$ 通过 $\sigma:G\times U\to U$ 作用于集合 $U$. 记 $[U/G]$ 为习题~\bookref{1.4.8}~中定义的商群胚: 对象是元素 $u\in U$, 态射 $u\to u'$ 是满足 $u'=gu$ 的元素 $g\in G$. 记投影为 $p:U\to[U/G]$.
\begin{enumerate}[(a)]
\item 设 $P$ 是带有 $G$ 的自由作用的集合, 其商为 $T=P/G$. 若 $f:P\to U$ 是 $G$-等变映射, 证明有 Cartesian 图
\[
\begin{tikzcd}
P\arrow[r,"f"]\arrow[d]\arrow[dr,phantom,"\square"]&U\arrow[d,"p"]\\
T\arrow[r]&{[U/G]},
\end{tikzcd}
\]
即 $P$ 等价于这个纤维积.
\item 证明有 Cartesian 图
\[
\begin{tikzcd}
G\times U\arrow[r,"\sigma"]\arrow[d,"p_2"']\arrow[dr,phantom,"\square"]&
U\arrow[d,"p"]\\
U\arrow[r,"p"']&{[U/G]}
\end{tikzcd}
\qquad\text{以及}\qquad
\begin{tikzcd}[column sep=large]
G\times U\arrow[r,"{(\sigma,p_2)}"]\arrow[d]\arrow[dr,phantom,"\square"]&
U\times U\arrow[d,"p\times p"]\\
{[U/G]}\arrow[r,"\Delta"']&{[U/G]\times[U/G]}.
\end{tikzcd}
\]
\end{enumerate}
\end{exercise}
\begin{exercise}[有益的练习]\label{src:3.4.31}
回顾例~\bookref{1.4.3}, 群 $G$ 的分类群胚 $BG$ 是只有一个对象 $*$ 且 $\Mor(*,*)=G$ 的范畴.
\begin{enumerate}[(a)]
\item 设 $\varphi:H\to G$ 是群同态. 证明有诱导的群胚态射 $BH\to BG$, 且 $BH\times_{BG}\mathrm{pt}\cong[G/H]$, 其中 $\mathrm{pt}$ 表示只有一个对象与一个态射的群胚.
\item 若 $K\trianglelefteq G$ 是正规子群, 其商为 $Q=G/K$, 证明有 Cartesian 图
\[
\begin{tikzcd}
Q\arrow[r]\arrow[d]\arrow[dr,phantom,"\square"]&
BK\arrow[r]\arrow[d]\arrow[dr,phantom,"\square"]&
\mathrm{pt}\arrow[d]\\
\mathrm{pt}\arrow[r]&BG\arrow[r]&BQ.
\end{tikzcd}
\]
\item 设群 $G$ 作用于集合 $U$, $[U/G]$ 是群胚商. 若点 $u\in U$ 的稳定子为 $G_u$, 轨道为 $Gu$, 证明有群胚态射 $BG_u\to[U/G]$ 以及 Cartesian 图
\[
\begin{tikzcd}
Gu\arrow[r]\arrow[d]\arrow[dr,phantom,"\square"]&U\arrow[d,"p"]\\
BG_u\arrow[r]&{[U/G]}.
\end{tikzcd}
\]
\end{enumerate}
\end{exercise}
\begin{exercise}[重要]\label{src:3.4.32}
证明群胚 $\mathcal C$ 等价于集合当且仅当 $\mathcal C\to\mathcal C\times\mathcal C$ 完全忠实.
\end{exercise}

\subsection{预叠的纤维积}\label{src:3.4.7-subsection}
范畴 $\mathcal S$ 上预叠的态射 $\mathcal X\to\mathcal Y$ 与 $\mathcal Y'\to\mathcal Y$ 的纤维积, 是预叠 $\mathcal X\times_{\mathcal Y}\mathcal Y'$, 其在 $S\in\mathcal S$ 上的纤维范畴为群胚的纤维积 $\mathcal X(S)\times_{\mathcal Y(S)}\mathcal Y'(S)$.
\begin{construction}\label{src:3.4.33}
设 $f:\mathcal X\to\mathcal Y$ 与 $g:\mathcal Y'\to\mathcal Y$ 是范畴 $\mathcal S$ 上的预叠态射. 定义 $\mathcal S$ 上的预叠 $\mathcal X\times_{\mathcal Y}\mathcal Y'$ 为由三元组 $(x,y',\gamma)$ 组成的范畴, 其中对象 $x\in\mathcal X$, $y'\in\mathcal Y'$ 位于同一个对象 $S:=p_{\mathcal X}(x)=p_{\mathcal Y'}(y')\in\mathcal S$ 上, 而 $\gamma:f(x)\xrightarrow{\sim}g(y')$ 是 $\mathcal Y(S)$ 中的同构.
态射 $(x_1,y_1',\gamma_1)\to(x_2,y_2',\gamma_2)$ 由三元组 $(h,\chi,\lambda)$ 组成, 其中
\[
h:p_{\mathcal X}(x_1)=p_{\mathcal Y'}(y_1')
\longrightarrow p_{\mathcal Y'}(y_2')=p_{\mathcal X}(x_2)
\]
是 $\mathcal S$ 中的态射, 而 $\chi:x_1\to x_2$ 与 $\lambda:y_1'\to y_2'$ 分别是 $\mathcal X$ 与 $\mathcal Y'$ 中位于 $h$ 上的态射, 使得
\[
\begin{tikzcd}
f(x_1)\arrow[r,"f(\chi)"]\arrow[d,"\gamma_1"']&
f(x_2)\arrow[d,"\gamma_2"]\\
g(y_1')\arrow[r,"g(\lambda)"']&g(y_2')
\end{tikzcd}
\]
交换.
\end{construction}
记投影 $(x,y',\gamma)\mapsto x$ 与 $(x,y',\gamma)\mapsto y'$ 为 $p_1:\mathcal X\times_{\mathcal Y}\mathcal Y'\to\mathcal X$ 与 $p_2:\mathcal X\times_{\mathcal Y}\mathcal Y'\to\mathcal Y'$. 定义 $2$-同构 $\alpha:f\circ p_1\xrightarrow{\sim}g\circ p_2$, 其在对象 $(x,y',\gamma)\in\mathcal X\times_{\mathcal Y}\mathcal Y'$ 上定义为 $\alpha_{(x,y',\gamma)}:=\gamma:f(x)\to g(y')$. 由此得到 $2$-交换图
\begin{equation}\label{src:3.4.34}
\begin{tikzcd}[column sep=large,row sep=large]
\mathcal X\times_{\mathcal Y}\mathcal Y'\arrow[r,"p_2"]\arrow[d,"p_1"']&
\mathcal Y'\arrow[d,"g"]\\
\mathcal X\arrow[r,"f"']&\mathcal Y
\arrow[Rightarrow,from=2-1,to=1-2,shorten <=1.6em,shorten >=1.6em,"\alpha" description]
\end{tikzcd}
\end{equation}
\begin{theorem}\label{src:3.4.35}
预叠 $\mathcal X\times_{\mathcal Y}\mathcal Y'$ 连同~\bookeqref{3.4.34}~中的态射 $p_1,p_2$ 与 $2$-同构 $\alpha$ 满足以下泛性质: 对每个带有 $2$-同构 $\tau:f\circ q_1\xrightarrow{\sim}g\circ q_2$ 的 $2$-交换图
\[
\begin{tikzcd}[column sep=large,row sep=large]
{}&\mathcal Y'\arrow[dr,"g"]&{}\\
\mathcal T\arrow[ur,bend left=25,"q_2"]\arrow[dr,bend right=25,"q_1"']&
\mathcal X\times_{\mathcal Y}\mathcal Y'\arrow[u,"p_2"']\arrow[d,"p_1"]&
\mathcal Y\\
{}&\mathcal X\arrow[ur,"f"']&{}
\arrow[Rightarrow,from=3-1,to=1-1,shorten <=3.5em,shorten >=3.5em,shift right=3em,"\tau"']
\arrow[Rightarrow,from=3-3,to=1-3,shorten <=3.5em,shorten >=3.5em,shift left=3em,"\alpha"]
\end{tikzcd}
\]
存在态射 $h:\mathcal T\to\mathcal X\times_{\mathcal Y}\mathcal Y'$ 与 $2$-同构 $\beta:q_1\to p_1\circ h$, $\rho:q_2\to p_2\circ h$, 给出 $2$-交换图
\[
\begin{tikzcd}[column sep=large,row sep=large]
{}&\mathcal Y'\arrow[dr,"g"]&{}\\
\mathcal T\arrow[r,"h"]\arrow[ur,bend left=25,"q_2"]
\arrow[dr,bend right=25,"q_1"']&
\mathcal X\times_{\mathcal Y}\mathcal Y'\arrow[u,"p_2"']\arrow[d,"p_1"]&
\mathcal Y\\
{}&\mathcal X\arrow[ur,"f"']&{}
\arrow[Rightarrow,from=3-1,to=2-1,shorten <=1.2em,shorten >=1.2em,shift right=3em,"\beta"']
\arrow[Rightarrow,from=1-1,to=2-1,shorten <=1.2em,shorten >=1.2em,shift left=3em,"\rho"]
\arrow[Rightarrow,from=3-3,to=1-3,shorten <=3.5em,shorten >=3.5em,shift left=3em,"\alpha"]
\end{tikzcd}
\]
使得图
\[
\begin{tikzcd}[column sep=large]
f\circ q_1\arrow[r,"f(\beta)"]\arrow[d,"\tau"']&
f\circ p_1\circ h\arrow[d,"\alpha\circ h"]\\
g\circ q_2\arrow[r,"g(\rho)"']&g\circ p_2\circ h
\end{tikzcd}
\]
交换. 数据 $(h,\beta,\rho)$ 在唯一同构的意义下唯一.
\end{theorem}
\begin{proof}
在对象上, 将 $h:\mathcal T\to\mathcal X\times_{\mathcal Y}\mathcal Y'$ 定义为
\[
t\longmapsto\bigl(q_1(t),q_2(t),\tau_t:f(q_1(t))\xrightarrow{\sim}g(q_2(t))\bigr);
\]
在态射上, 定义为 $(\Psi:t\to t')\mapsto(p_{\mathcal T}(\Psi),q_1(\Psi),q_2(\Psi))$. 有函子等式 $q_1=p_1\circ h$ 与 $q_2=p_2\circ h$, 因而可以将 $\beta,\rho$ 定义为恒等自然变换. 其余细节留给读者.
\end{proof}
\begin{definition}\label{src:3.4.36}
我们称 $2$-交换图
\[
\begin{tikzcd}[column sep=large,row sep=large]
\mathcal X'\arrow[r]\arrow[d]&\mathcal Y'\arrow[d]\\
\mathcal X\arrow[r]&\mathcal Y
\arrow[Rightarrow,from=2-1,to=1-2,shorten <=1.4em,shorten >=1.4em,"\alpha" description]
\end{tikzcd}
\]
为 Cartesian 图, 若它满足定理~\bookref{3.4.35}~的泛性质. 我们常将预叠的 Cartesian 图写成
\[
\begin{tikzcd}
\mathcal X'\arrow[r]\arrow[d]\arrow[dr,phantom,"\square"]&
\mathcal Y'\arrow[d]\\
\mathcal X\arrow[r]&\mathcal Y,
\end{tikzcd}
\]
其中 $2$-同构 $\alpha$ 的存在是隐含的.
\end{definition}
当 $\mathcal X,\mathcal Y$ 是 $\Sch$ 上的预叠时, 有 $\mathcal X\times_{\Sch}\mathcal Y=\mathcal X\times_{\Spec\mathbb Z}\mathcal Y$, 有时简记为 $\mathcal X\times\mathcal Y$. 类似地, 在域 $k$ 上工作时, 将 $\mathcal X\times_{\Sch/k}\mathcal Y=\mathcal X\times_{\Spec k}\mathcal Y$ 简记为 $\mathcal X\times_k\mathcal Y$, 或在不会混淆时记为 $\mathcal X\times\mathcal Y$.

\subsection{例子}\label{src:3.4.8-subsection}
下面的习题对于验证商叠的代数性及建立其性质至关重要 (例如见定理~\bookref{4.1.10}).
\begin{exercise}[重要, 有益的练习]\label{src:3.4.37}
设 $G\to S$ 是通过 $\sigma:G\times_S U\to U$ 作用于 $S$ 上概形 $U$ 的光滑仿射群概形, $[U/G]$ 是商叠 (定义~\bookref{3.4.15}).
\begin{enumerate}[(a)]
\item 设态射 $T\to[U/G]$ 通过 $2$-Yoneda 引理~(\bookref{3.4.21})~对应于主 $G$-丛 $P\to T$ 与 $G$-等变映射 $f:P\to U$. 证明有 Cartesian 图
\[
\begin{tikzcd}
P\arrow[r,"f"]\arrow[d]\arrow[dr,phantom,"\square"]&U\arrow[d]\\
T\arrow[r]&{[U/G]}.
\end{tikzcd}
\]
(后面将看到, $[U/G]$ 是代数叠且 $U\to[U/G]$ 是主 $G$-丛 (定理~\bookref{4.1.10}). 这使我们能够将主 $G$-丛 $U\to[U/G]$ 连同恒等映射 $U\to U$, 等同于 $[U/G]$ 上的万有族, 它通过 $2$-Yoneda 引理对应于恒等映射 $[U/G]\to[U/G]$.)
\item 证明有 Cartesian 图
\[
\begin{tikzcd}
G\times_S U\arrow[r,"\sigma"]\arrow[d,"p_2"']\arrow[dr,phantom,"\square"]&
U\arrow[d,"p"]\\
U\arrow[r,"p"']&{[U/G]}
\end{tikzcd}
\qquad\text{以及}\qquad
\begin{tikzcd}[column sep=large]
G\times_S U\arrow[r,"{(\sigma,p_2)}"]\arrow[d]\arrow[dr,phantom,"\square"]&
U\times_S U\arrow[d,"p\times p"]\\
{[U/G]}\arrow[r,"\Delta"']&{[U/G]\times_S[U/G]}.
\end{tikzcd}
\]
\end{enumerate}
\end{exercise}
下一道习题中的图, 如同在概形情形中一样, 将得到广泛使用. 我们将在 \S\bookref{4.2.2-subsection} 中用它定义稳定子与惯性叠.
\begin{exercise}[重要]\label{src:3.4.38}
设 $\mathcal X$ 是范畴 $\mathcal S$ 上的预叠. 设 $S,T\in\mathcal S$ 是对象, 通过例~\bookref{3.4.9}~将其视为 $\mathcal S$ 上的预叠. 证明对所有态射 $a:S\to\mathcal X$ 与 $b:T\to\mathcal X$, 都有 Cartesian 图
\[
\begin{tikzcd}[column sep=large]
S\times_{\mathcal X}T\arrow[r]\arrow[d]\arrow[dr,phantom,"\square"]&
S\times T\arrow[d,"a\times b"]\\
\mathcal X\arrow[r,"\Delta"']&\mathcal X\times\mathcal X,
\end{tikzcd}
\]
其中纤维积 $S\times T$ 与 $\mathcal X\times\mathcal X$ 都在 $\mathcal S$ 上取.
\end{exercise}
\begin{exercise}[Isom 预层, 重要]\label{src:3.4.39}
设 $\mathcal X$ 是范畴 $\mathcal S$ 上的预叠. 对 $S\in\mathcal S$, 回顾例~\bookref{3.2.12}, 限制范畴 $\mathcal S/S$ 的对象是 $\mathcal S$ 中的态射 $T\to S$, 态射是 $S$-态射.
\begin{enumerate}[(a)]
\item 证明对 $\mathcal X$ 中位于 $S$ 上的对象 $a,b$, 函子
\[
\begin{aligned}
\Isom_{\mathcal X(S)}(a,b):\mathcal S/S&\longrightarrow\Sets,\\
(T\xrightarrow{f}S)&\longmapsto\Mor_{\mathcal X(T)}(f^*a,f^*b),
\end{aligned}
\]
定义 $\mathcal S/S$ 上的预层, 其中 $f^*a,f^*b$ 是选定的拉回.
\item 证明有 Cartesian 图
\[
\begin{tikzcd}[column sep=large]
\Isom_{\mathcal X(S)}(a,b)\arrow[r]\arrow[d]\arrow[dr,phantom,"\square"]&
S\arrow[d,"{(a,b)}"]\\
\mathcal X\arrow[r,"\Delta"']&\mathcal X\times\mathcal X.
\end{tikzcd}
\]
\item 证明预层 $\Aut_{\mathcal X(T)}(a)=\Isom_{\mathcal X(T)}(a,a)$ 自然是群值预层.
\item 证明 $\mathcal X$ 等价于预层当且仅当对角态射 $\mathcal X\to\mathcal X\times\mathcal X$ 完全忠实.
\end{enumerate}
\end{exercise}
\begin{exercise}[有益的练习]\label{src:3.4.40}
\begin{enumerate}[(a)]
\item 若 $H\to G$ 是概形 $S$ 上的光滑仿射群概形态射, 利用定义~\bookref{B.1.62}~从主 $H$-丛构造主 $G$-丛, 定义 $\Sch/S$ 上的预叠态射 $BH\to BG$.
\item 证明 $BH\times_{BG}S\cong[G/H]$.
\item 若 $1\to K\to G\to Q\to1$ 是域 $k$ 上光滑仿射代数群的正合序列, 证明有 Cartesian 图
\[
\begin{tikzcd}
Q\arrow[r]\arrow[d]\arrow[dr,phantom,"\square"]&
BK\arrow[r]\arrow[d]\arrow[dr,phantom,"\square"]&
\Spec k\arrow[d]\\
\Spec k\arrow[r]&BG\arrow[r]&BQ.
\end{tikzcd}
\]
\end{enumerate}
\end{exercise}
\begin{exercise}[有益的练习]\label{src:3.4.41}
设 $G,H$ 是概形 $S$ 上的光滑仿射群概形.
\begin{enumerate}[(a)]
\item 证明 $B(G\times_S H)\cong BG\times_S BH$.
\item 若 $X,Y$ 是分别带有 $G,H$ 作用的 $S$-概形, 证明 $[(X\times_S Y)/(G\times_S H)]\cong[X/G]\times_S[Y/H]$.
\item 推出在 $\Spec\mathbb Z$ 上有
\[
[\mathbb A^n/\mathbb G_m^n]\cong
\underbrace{[\mathbb A^1/\mathbb G_m]\times\cdots\times[\mathbb A^1/\mathbb G_m]}_{n\text{ 次}}.
\]
\end{enumerate}
\end{exercise}
\begin{exercise}[重要, 有益的练习]\label{src:3.4.42}
仿照光滑曲线的预叠 $\mathcal M_g$ (例~\bookref{3.4.10}), 设 $\mathcal M_{g,1}$ 是如下预叠: 概形 $S$ 上的对象为光滑曲线族 $C\to S$ 及截面 $\sigma:S\to C$. 记遗忘截面的预叠态射为 $\mathcal M_{g,1}\to\mathcal M_g$. 证明, 若态射 $S\to\mathcal M_g$ 对应于曲线族 $C\to S$, 则有 Cartesian 图
\[
\begin{tikzcd}
C\arrow[r]\arrow[d]\arrow[dr,phantom,"\square"]&\mathcal M_{g,1}\arrow[d]\\
S\arrow[r]&\mathcal M_g.
\end{tikzcd}
\]
换言之, $\mathcal M_{g,1}\to\mathcal M_g$ 是万有族, 其定义将在定义~\bookref{4.1.29}~中给出.
\end{exercise}
\begin{exercise}[有益的练习]\label{src:3.4.43}
设 $H,G$ 是概形 $S$ 上的光滑仿射群概形. 记 $\Hom(H,G)$ 为 $(\Sch/S)_{\text{ét}}$ 上的层, 其在 $S$-概形 $T$ 上的截面是同态 $H_T\to G_T$. 记 $\Mor(BH,BG)$ 为 $S$ 上的预叠, 其位于 $T\in S$ 上的对象是态射 $B(H_T)\to B(G_T)$. 证明
\[
\Mor(BH,BG)\cong[\Hom(H,G)/G],
\]
其中 $G$ 通过共轭作用.
\end{exercise}
